\naslov{Nesimetrični problem lastnih vrednosti}

Naj bosta $A, E \in \R^{n \times n}$ matriki in $\varepsilon \in \R$.
Označimo lastne vrednosti matrike $A + \varepsilon E$ z
$\lambda_i(\varepsilon)$.

\begin{izrek}[Bauer-Fike]
  Naj bo $A = X D X^{-1}$ diagonalizabilna z $D = \operatorname{diag}(\lambda_1,
  \ldots, \lambda_n)$ in $X = [x_1, \ldots, x_n]$.
  Za vsak $\varepsilon > 0$ lastne vrednosti $A + \varepsilon E$ ležijo v uniji
  krogov
  \[
	K_i = \{ z \such \abs{z - \lambda_i} \le \varepsilon \norm{E} \norm{X}
	\norm{X^{-1}} \},
  \]
  kjer je $\norm{\cdot}$ $1$-, $\infty$- ali $2$-norma.
  Če unija razpade na več povezanih komponent, vsaka vsebuje toliko lastnih
  vrednosti, kolikor krogov jo sestavlja.
\end{izrek}

\begin{proof}
  Naj bo $\lambda(\varepsilon)$ lastna vrednost $A + \varepsilon E$.
  Matrika $A + \varepsilon E - \lambda(\varepsilon) I$ je singularna, velja
  \[
	A + \varepsilon E - \lambda(\varepsilon) I
	= X(D + \varepsilon X^{-1} E X - \lambda(\varepsilon) I) X^{-1}.
  \]
  Če je $\lambda(\varepsilon) = \lambda_i$ za nek $i$, očitno leži v uniji.
  Sicer je $D - \lambda(\varepsilon) I$ nesingularna, ker pa je srednji člen
  produkta zgoraj singularen, je singularna tudi
  \[
	I + \varepsilon (D - \lambda(\varepsilon) I)^{-1} X^{-1} E X.
  \]
  Torej je
  \[
	1 \le \norm{\varepsilon (D - \lambda(\varepsilon) I)^{-1} X^{-1} E X}
	\le \varepsilon \norm{(D - \lambda(\varepsilon) I)^{-1}} \norm{X^{-1}}
	\norm{X} \norm{E}.
  \]
  Elementi matrike $(D - \lambda(\varepsilon))^{-1}$ so oblike $(\lambda_i -
  \lambda(\varepsilon))^{-1}$, torej je res
  \[
	\min_i \abs{\lambda_i - \lambda(\varepsilon)} \le \varepsilon \norm{X^{-1}}
	\norm{E} \norm{X}.
  \]
\end{proof}

\vprasanje{Povej in dokaži Bauer-Fikeov izrek o občutljivosti problema lastnih
  vrednosti.}

\begin{posledica}
  Če je $A$ simetrična, potem za vsako lastno vrednost $\tilde{\lambda}$ matrike
  $A + \varepsilon E$ obstaja lastna vrednost $\lambda_i$ matrike $A$, da je
  $\abs{\tilde{\lambda} - \lambda_i} \le \varepsilon \norm{E}_2$.
\end{posledica}

\begin{proof}
  Matriko lahko diagonaliziramo z ortogonalno bazo, torej sta normi $\norm{X}_2
  = \norm{X^{-1}}_2 = 1$.
\end{proof}

\vprasanje{Kaj pravi Bauer-Fikeov izrek o simetričnih matrikah?}

Naj bo $Ax = \lambda x$ in $(A + \Delta A)(x + \Delta x) = (\lambda + \Delta
\lambda)(x + \Delta x)$.
Recimo, da je $\lambda$ enostavna lastna vrednost in $y$ levi lastni vektor,
$y^H A = \lambda y^H$.
Potem je
\[
  Ax + \Delta A x + A \Delta x + \Delta A \Delta x = \lambda x + \Delta \lambda
  x + \lambda \Delta x + \Delta \lambda \Delta x.
\]
Zadnji člen na vsaki strani zanemarimo in množimo z leve z $y^H$, da dobimo
\[
  y^H \Delta A x + y^H A \Delta x = \Delta \lambda y^H x + \lambda y^H \Delta x,
\]
oziroma
\[
  \Delta \lambda = \frac{y^H \Delta A x}{y^H x}.
\]

\vprasanje{Izpelji predpis za motnjo v lastni vrednosti ob dani motnji matrike
  $\Delta A$.}

\begin{izrek}
  Naj bo $\lambda_i$ enostavna lastna vrednost $A$ z normiranima desnim in levim
  lastnim vektorjem $x_i$ in $y_i$.
  Če je $\lambda_i(\varepsilon)$ ustrezna lastna vrednost $A + \varepsilon E$,
  potem velja
  \[
	\lambda_i(\varepsilon) = \lambda_i + \varepsilon \frac{y_i^H E x_i}{y_i^H
	  x_i} + O(\varepsilon^2).
  \]
\end{izrek}

Definiramo
\[
  s_i = \frac{y_i^H x_i}{\norm{y_i}_2 \norm{x_i}_2}.
\]
Potem je občutljivost lastne vrednosti enaka $1/\abs{s_i}$.
Ocenimo lahko $\abs{\lambda_i(\varepsilon) - \lambda_i} \le \varepsilon
\frac{\norm{E}_2}{\abs{s_i}} + O(\varepsilon^2)$.
Do maksimalne spremembe pride pri $E = y_i x_i^H$; tedaj je $\abs{y_i^H E x_i} =
1$.

\vprasanje{Kaj je občutljivost lastne vrednosti?}

Če je $\lambda$ $p$-kratna lastna vrednost in ima v Jordanovi formi kletke
velikosti $m_1, \ldots, m_k$ za $m_1 + \cdots + m_k = p$, lahko pričakujemo
$\abs{\lambda_i(\varepsilon) - \lambda_i} = O(\varepsilon^{1/m_1})$ za ureditev
$m_1 \ge m_2 \ge \cdots \ge m_k$.

\begin{izrek}
  Naj bo $A = X D X^{-1}$ diagonalizabilna ter $Y^H A = D Y^H$, kjer je $X =
  [x_1, \ldots, x_n]$, $Y = [y_1, \ldots, y_n]$ in $\norm{x_i}_2 = \norm{y_i}_2
  = 1$.
  Če je $\lambda_i$ enostavna lastna vrednost, potem za dovolj majhen
  $\varepsilon$ za lastni vektor $x_i(\varepsilon)$ za $\lambda_i(\varepsilon)$
  velja
  \[
	x_i(\varepsilon) = x_i + \sum_{i \ne j} \varepsilon \frac{y_j^H E
	  x_i}{(\lambda_i - \lambda_j) s_j} x_j + O(\varepsilon^2).
  \]
\end{izrek}

\begin{izrek}
  Če je $\lambda_i$ enostavna lastna vrednost $A$, je $s_i \ne 0$.
\end{izrek}

\begin{proof}
  Recimo, da je $y_i^H x_i = 0$.
  Naj bo $U$ taka unitarna matrika, da je $U e_1 = x_i$.
  Oglejmo si
  \[
	B = U^H A U =
	\begin{bmatrix}
	  \lambda_i & \cdots \\
	  0 & C
	\end{bmatrix}.
  \]
  Velja $B e_1 = \lambda_i e_1$, levi lastni vektor je oblike $z^H = y_i^H U$.
  Vemo $y_i^H x_i = 0$, torej je $y_i^H U U^H x_i = 0$, ker pa je $U^H x_i =
  e_1$, velja $y_i^H U = [0, w^H]$.
  Torej $[0, w^H] B = \lambda_i [0, w^H]$, oziroma $w^H C = \lambda_i w^H$.
  Torej je $\lambda_i$ vsaj dvojna lastna vrednost $A$.
\end{proof}

\vprasanje{Pokaži, da je občutljivost enostavne lastne vrednosti dobro
  definirana.}

\podnaslov{Implicitna QR iteracija}

\begin{izrek}[Izrek o implicitnem Q]
  Naj bo $Q = [q_1, \ldots, q_n]$ taka ortogonalna matrika, da je $H = Q^T A Q$
  nerazcepna zgornje Hessenbergova matrika.
  Potem so stolpci $q_2, \ldots, q_n$ do predznaka natančno določeni s $q_1$.
\end{izrek}

\begin{proof}
  Naj bo tudi $V^T A V = G$, kjer je $G$ nerazcepna zgornje Hessenbergova, $V$
  ortogonalna in $v_1 = q_1$.
  Velja $AV = VG$ in $Q^T A = H Q^T$.
  Če prvo enačbo množimo z leve s $Q^T$ in drugo z desne z $V$, dobimo $Q^T V G
  = Q^T A V = H Q^T V$, za $W = Q^T V$ torej velja $W G = H W$.
  Matrika $W$ je ortogonalna z $w_1 = e_1$, ker je $q_1 = v_1$.
  Poglejmo sedaj $i$-ti stolpec matrike $WG = HW$.
  Velja
  \[
	\sum_{j=1}^{i+1} g_{ji} w_j = H w_i,
  \]
  torej
  \[
	g_{i+1,i} w_{i+1} = H w_i - \sum_{j=1}^i g_{ji} w_j.
  \]
  Pokažimo, da ima $w_i$ neničelnih le prvih $i$ elementov.
  Indukcija na $i$;
  \[
	w_{i+1} = \inv{g_{i+1,i}} H w_i - \inv{g_{i+1,i}} \sum_{j=1}^i g_{ji} w_j.
  \]
  Matrika $H$ je zgornje Hessenbergova, drug člen razlike pa po indukcijski
  predpostavki vsebuje le elemente na mestih $1$ do $i$.
  Torej je $W$ zgornje trikotna; ker je tudi ortogonalna, mora biti diagonalna z
  $\pm 1$ na diagonali.
\end{proof}

\vprasanje{Povej in dokaži izrek o implicitnem Q.}

Obravnavajmo QR iteracijo z enojnim premikom.
Če začnemo z zgornje Hessenbergovo matriko $A_0$, v vsakem koraku izračunamo QR
razcep $A_k - \sigma_k I = Q_k R_k$, in nato določimo $A_{k+1} = R_k Q_k +
\sigma_k I$.
Pišemo lahko tudi $A_{k+1} = Q_k^T A_k Q_k$.
Ker smo začeli z zgornje Hessenbergovo matriko, naredimo QR razcep z
Givensovimi rotacijami.
Prvo z leve množimo z rotacijo $R_{12}^T$, ki nam spremeni prvi vrstici matrike;
pri množenju z desne z $R_{12}$ pa dobimo nov element na mestu $(3,1)$.
V nadaljevanju bomo množili z novimi rotacijami tako, da ne bomo spremenili
prvega stolpca v prehodni matriki na desni, nov element pa bomo počasi premikali
navzdol izven matrike.

\vprasanje{Kako uporabiš izrek o implicitnem Q pri QR iteraciji z enojnim
  premikom?}

Za dvojni premik velja
\[
  N_k = A_k^2 - (\sigma_{k1} + \sigma_{k2}) A_k + \sigma_{k1} \sigma_{k2} I =
  Q_k R_k
\]
in $A_{k+1} = Q_k^T A_k Q_k$.
Prvi stolpec matrike $N_k$ lahko izračunamo v konstantno mnogo operacijah.
Iz tega določimo Householderjevo zrcaljenje $P_1$, ki stolpec zrcali v $e_1$.
Tukaj dobimo nove elemente tako pri množenju z leve kot pri množenju z desne,
tako da premikamo grbo velikosti 2x2.

\vprasanje{Kako uporabiš izrek o implicitnem Q pri QR iteraciji z dvojnim
  premikom?}

% LocalWords:  neničelnih Givensovimi Householderjevo
